\scanpage{133}
\deleted{\unsure{sous-groupe}} isomorphe à $\mathbb U=\{z\in\mathbb C\mid |z|=1\}$ tel que l'action sur $V_{\mathbb C}$ ait exactement 2 sous-espaces propres,
% slanted
\margin{(c'est automatique si $\mathrm{rg}_{\mathbb R}V=2$)}

$2^\circ$) D'un des deux éléments d'ordre $4$ de $\mathbb U$, ou ce qui revient au $=$, d'une orientation de $\mathbb U$, ou encore \inserted{(si $\mathrm{rg}_{\mathbb R}V=2$, plus général\textsuperscript{t} si $\mathrm{rg}_{\mathbb R}V=2n$, $n$ impair)} d'une orientation de $V$
% slanted
\margin{ou encore d'un \inserted{des deux} isomorphismes $\mathbb U\xrightarrow{\sim}U$, ou $U\xrightarrow{\sim}\mathbb U$}

-- l'élément d'ordre $4$ de $\mathbb U$ associé à une orientation $\omega$ de \inserted{$U$ (ou de} $V$) sera noté $u_\omega$, c'est donc la rotation de $+\pi/2$ pour l'orientation \inserted{$\omega$} de $V$ et \unsure{leur|la|une} structure \deleted{complexe}\inserted{euclidienne} sur $V$ choisie invariante sous $U$ (\uline{NB} \inserted{si $\mathrm{rg}_{\mathbb R}V=2$,} une telle \deleted{struct}\inserted{métrique} est unique à scalaire multiplicatif près).

Alors $V_{\mathbb C}$ sous l'action de $U$ se décompose en deux sous-espaces propres
\[
 V_{\mathbb C}=\bigoplus_{\omega\in\underline\omega}\deleted{\unsure{\omega}}L_\omega \tag*{(\unsure{14|13})}
\]
où pour $\omega\in\underline\omega(U)=\underline\omega$, $L_\omega$ est caractérisé par le fait que l'isomorphisme
\[
 \varphi_\omega\colon U\xrightarrow{\sim}\mathbb U\qquad g(z)=\varphi_\omega(g)\cdot z\quad\text{pour }z\in L_\omega \tag*{(15)}
\]
défini par l'action de $U$ sur $L_\omega$, corresponde à l'orientation $\omega$. La donnée d'une structure complexe sur $V$ \inserted{compatible avec $U$} équivaut alors à celle d'un des deux isomorphismes $\varphi_\omega$ (permettant de faire opérer $\mathbb U$ sur $V$, donc aussi $\mathbb R^{\times+}\times\mathbb U\simeq\mathbb C^\times$, \deleted{\unsure{donc|dans}} $\mathbb C$ \ldots). Ceci dit, pour \deleted{$\unsure{x\in\mathbb C|z\in\mathbb C}$} $z\in L_\omega$, on aura
\[
 \begin{aligned}
 \Re (\varphi_\omega(g)z)=\Re (gz)=g\Re (z)&=\text{produit de}\\
 &\Re (z)\ \text{par }\varphi_\omega(g)\ \text{pour la structure complexe}\\[-2pt]
 &\text{de }V\text{ associée à }\varphi_\omega
 \end{aligned}
\]
% archive sheet number: 132
